% TOP_PROBLEM: {"id": 30006850, "title": "Unrestricted Linear-Circuit Lower Bound for the Discrete Fourier Transform", "category_id": 15, "category": {"id": 15, "name": "computer_science", "display_name": "Computer Science", "description": "Computational complexity, algorithms, and theoretical CS.", "slug": "computer-science", "order_index": 15, "created_at": "2026-07-31T15:26:25.671Z"}, "status": "open"}
% FIRST_POSTED: 2026-09-25
% ENTRY_KIND: statement_only
% !TeX program = lualatex
% Definition and sources only; this file is not an LLM solution attempt.
\documentclass[11pt]{article}
\usepackage[margin=25mm]{geometry}
\usepackage{fontspec}
\setmainfont{Latin Modern Roman}
\usepackage{amsmath,amssymb,mathtools}
\usepackage{unicode-math}
\setmathfont{Latin Modern Math}
\usepackage{xurl,hyperref}
\hypersetup{colorlinks=true,urlcolor=blue}
\setcounter{secnumdepth}{0}
\setlength{\parindent}{0pt}
\setlength{\parskip}{5pt}
\setlength{\emergencystretch}{3em}
\begin{document}
\section*{\texorpdfstring{Unrestricted Linear-Circuit Lower Bound for the Discrete Fourier Transform}{Unrestricted Linear-Circuit Lower Bound for the Discrete Fourier Transform}}\label{30006850}
\subsection{Definitions and mathematical statement}
For $n\ge2$, let $\omega_n=e^{-2\pi i/n}$ and define the discrete Fourier transform by
\[
 y_j=\sum_{\ell=0}^{n-1}\omega_n^{j\ell}x_\ell,
 \qquad0\le j<n.
\]
A linear straight-line circuit starts with the input coordinates and allowed scalar constants. Each gate forms $\lambda f+\mu g$ from two previously available values, with arbitrary complex coefficients; its cost is one gate. Let $C_{\rm DFT}(n)$ be the minimum gate count for producing all outputs exactly. The conjecture is $C_{\rm DFT}(n)\ge c n\log n$ for an absolute $c>0$ and all sufficiently large $n$. There is no bound on coefficient magnitude or intermediate conditioning. Consequently, a lower bound proved only for bounded-coefficient or well-conditioned circuits is an approach to, not a solution of, the unrestricted question.
\par\smallskip\textit{Formulation and concept references:} \href{https://arxiv.org/pdf/1403.1307v5}{[S1]}.

\subsection{Short English statement}
Must every exact linear circuit for the discrete Fourier transform use order $n\log n$ operations, even with arbitrary complex coefficients and ill-conditioned intermediate calculations?
\subsection{Sources}
\begin{itemize}
\item[S1]An Omega((n log n)/R) Lower Bound for Fourier Transform Computation in the R-Well Conditioned Model, v5 header24July2014; internal31October2018.
\url{https://arxiv.org/pdf/1403.1307v5}.
\textit{Formulation source.}
\end{itemize}
\end{document}
